\section{Deutsch Algorithm}\label{sec:deutsch-algorithm}
The Deutsch Algorithm is one of the first quantum algorithms which showed more efficient than their classical counterpart.

Let $f:\B\longrightarrow\B$ be a function.
There are only four distinct functions, the constant functions $0$ and $1$, and the balanced functions identity and the opposite or flip function.
If we need to know if $f$ is constant or balanced, we need to evaluate $f$ twice, but with a quantum algorithm we only need evaluate once, as was shown by Deutsch~\cite{Deutsch1985}.

The key observation is that $f$ is constant if and only if $f(0)\oplus f(1)=0$.
The quantum algorithm extracts exactly this single bit of information with one evaluation of $f$, by querying the oracle on a superposition of both inputs.

Let $\gate{U}$ be the unitary operator defined on the computational basis by
\[
	\ket{xy}\gate{U}=\ket{x(y\oplus f(x))},
\]
where $\oplus$ is the addition mod 2.
If we apply the operator $\gate{U}$ over the $\ket{+-}$ state we have:
\[
	\begin{split}
	\ket{+-}\gate{U} &= \frac{1}{\sqrt{2}}(\ket{+0}-\ket{+1})\gate{U}=\frac{1}{2}(\ket{00}+\ket{10}-\ket{01}-\ket{11})\gate{U}=\\
	&= \frac{1}{2}(\ket{0(0\oplus f(0))}+\ket{1(0\oplus f(1))}-\ket{0(1\oplus f(0))}-\ket{1(1\oplus f(1))})=\\
		&=\begin{cases}
				\frac{1}{2}(\ket{00}+\ket{10}-\ket{01}-\ket{11})=\ket{+-}\text{ if $f(0)=0$ and $f(1) = 0$}\\
				\frac{1}{2}(\ket{00}+\ket{11}-\ket{01}-\ket{10})=\ket{--}\text{ if $f(0)=0$ and $f(1) = 1$}\\
				\frac{1}{2}(\ket{01}+\ket{10}-\ket{00}-\ket{11})=-\ket{--}\text{ if $f(0)=1$ and $f(1) = 0$}\\
				\frac{1}{2}(\ket{01}+\ket{11}-\ket{00}-\ket{10})=-\ket{+-}\text{ if $f(0)=1$ and $f(1) = 1$}
	\end{cases}
	\end{split}
\]
so, the resulting state is, up to a global phase:
\[
	\ket{+-}\gate{U}=(-1)^{f(0)}\begin{cases}
		            \ket{+-} & \text{ if $f$ is constant} \\
		            \ket{--} & \text{ if $f$ is balanced}
	\end{cases}.
\]

Let us consider a quantum circuit which implements the Deutsch algorithm.
According to the postfix convention, the circuit is read from left to right:
\[
	\gate{D}=\gate{X}_2\gate{H}_{1,2}\gate{U}\gate{H}_1
\]

Applying the circuit to the initial state $\ket{00}$ and using the computation above:
\begin{equation}
	\label{eq:deutchs-circuit}
	\begin{split}
		\ket{00}\gate{D}&=\ket{00}\gate{X}_2\gate{H}_{1,2}\gate{U}\gate{H}_1=\ket{01}\gate{H}_{1,2}\gate{U}\gate{H}_1=\ket{+-}\gate{U}\gate{H}_1=\\
		&=(-1)^{f(0)}
		\begin{cases}
			\ket{+-}\gate{H}_1\\
			\ket{--}\gate{H}_1
		\end{cases}
		=(-1)^{f(0)}
		\begin{cases}
			\ket{0-} & \text{ if $f$ is constant}\\
			\ket{1-} & \text{ if $f$ is balanced}
		\end{cases}
	\end{split}
\end{equation}

Therefore, measuring the first qubit in the computational basis yields $0$ if $f$ is constant and $1$ if $f$ is balanced, with a single evaluation of $\gate{U}$.
